% Version 0.25. Complete all-prime arithmetic-group torsion supplier.
\section{A complete torsion bound for the arithmetic homology group}
\label{sec:torsion-homology}

\status{TB1--TB10 prove the literal arithmetic-group supplier
\texttt{TorsBound}$(49)$ and the explicit stronger specialization
\texttt{TorsBoundPow}$(22)$.  The proof includes torsion at every prime.
It uses the source's ordered perfect-form nerve, but replaces the long
chain induction by a double-complex argument, the matrix normal-form
count by a row-module count, and the coordinate estimate by an inverse
Hermitian-matrix estimate.  No project-specific torsion, orbit-count,
acyclicity or stabilizer predicate is retained as an input.}

Write $T(A)$ for the torsion subgroup of an abelian group $A$.  In this
section homology has integral trivial coefficients unless specified.
All logarithms are natural.  The elementary background consists of
free resolutions, Smith normal form, finite-dimensional linear algebra,
and the integral lattice of a number field (for the latter, see
Milne, \href{https://www.jmilne.org/math/CourseNotes/ANT.pdf}{\emph{Algebraic
Number Theory}, Proposition~4.26}).  The exact constructions
that use this background are given below; in particular the acyclic
complex and the estimates entering the final bound are constructed.

\begin{proposition}[TB1: two integral torsion estimates]
\label{prop:tb-elementary-torsion}
Let $D:\mathbb Z^{(J)}\to\mathbb Z^a$ have columns of Euclidean norm
at most $b\ge1$, with $a<\infty$ and $J$ possibly infinite.  Then
$T(\operatorname{coker}D)$ is finite and has order at most $b^a$.
If $S$ is a finite group of order $s\ge1$, then, for every $q>0$,
\[
       |H_q(S,\mathbb Z)|\le s^{s^q}.
       \tag{TB.1}
\]
\end{proposition}

\begin{proof}
Let $L=\operatorname{im}D\subset\mathbb Z^a$ and let
$\overline L=(L\otimes\mathbb Q)\cap\mathbb Z^a$ be its saturation.
If its rank is zero the assertion is immediate.  Otherwise choose
$r=\operatorname{rank}L$ independent columns and let $L_0$ be their
span.  Choose $r$ coordinates on which these columns have nonzero
determinant.  Projection to these coordinates is injective on
$L\otimes\mathbb Q$, and consequently
\[
 [\overline L:L]\le[\overline L:L_0]
  =[\pi\overline L:\pi L_0]
  \le[\mathbb Z^r:\pi L_0]
  =|\det(\pi D_0)|\le b^r\le b^a.
\]
The penultimate bound is Hadamard's inequality, since coordinate
projection cannot increase a column norm.  The group
$\overline L/L$ is precisely the torsion of the cokernel.  This
also proves finiteness and does not require a bound on $|J|$.

For the second assertion use the inhomogeneous bar complex for $S$.
Its group in degree $q$ is free of rank $s^q$.  Its cycle subgroup is
therefore generated by at most $s^q$ elements, and so is its homology.
For completeness, the usual transfer to the trivial subgroup gives
\[
 H_q(S,\mathbb Z)\longrightarrow H_q(1,\mathbb Z)
       \longrightarrow H_q(S,\mathbb Z),
\]
whose composite is multiplication by $s$.  One constructs the first
map by summing all lifts of a simplex in the $s$-sheeted cover
$ES\to BS$ given by the bar construction; the boundary commutes
with this sum.  Projection sends each of those $s$ lifts back to
the original simplex, so the composite is multiplication by $s$.
For $q>0$ the intermediate homology is zero.  Hence $s$ annihilates
$H_q(S,\mathbb Z)$.  A group generated by $u$ elements and annihilated
by $s$ is a quotient of $(\mathbb Z/s)^u$, proving (TB.1), including
$s=1$.
\end{proof}

\begin{proposition}[TB2: torsion from an exact permutation complex]
\label{prop:tb-equivariant-bound}
Suppose a group $G$ acts on sets $X_p$, $p\ge0$, with equivariant
face maps deleting one of $p+1$ entries.  Suppose the augmented
complex
\[
 \cdots\longrightarrow\mathbb Z[X_1]\longrightarrow
 \mathbb Z[X_0]\longrightarrow\mathbb Z\longrightarrow0
 \tag{TB.2}
\]
is exact.  Fix $m\ge1$.  If there are at most $A$ orbits in each
$X_p$, $0\le p\le m$, and every stabilizer in those degrees is
finite of order at most $s\ge1$, then
\[
 |T(H_m(G,\mathbb Z))|
       \le (m+2)^A\,s^{m A s^m}.
 \tag{TB.3}
\]
In particular this torsion group is finite.
\end{proposition}

\begin{proof}
Take a free right $\mathbb ZG$-resolution $P_\bullet$ of $\mathbb Z$
and form the first-quadrant double complex
$P_q\otimes_{\mathbb ZG}\mathbb Z[X_p]$.  Exactness of (TB.2)
and freeness of each $P_q$ identify its total homology with
$H_*(G,\mathbb Z)$.  Computing the other direction first gives
\[
 E^1_{p,q}=\bigoplus_{[x]\in G\backslash X_p}
                    H_q(G_x,\mathbb Z)
       \ \Longrightarrow\ H_{p+q}(G,\mathbb Z).
 \tag{TB.4}
\]
Indeed, the permutation module of an orbit is
$\mathbb ZG\otimes_{\mathbb ZG_x}\mathbb Z$; restricting the free
resolution to $G_x$ proves the displayed identification directly.
Filtering the total complex by $p$ produces (TB.4); in each total
degree its filtration is finite, so the abutment has no convergence
qualification.

For $q>0$ and $p\le m$, TB1 bounds the corresponding group by
$s^{A s^q}$.  Its later-page subquotients have no larger order.
The row $q=0$ is the coinvariant complex
$\mathbb Z[G\backslash X_p]$.  In degree $m$ it has rank at most
$A$.  Every column of its incoming differential is a signed sum
of $m+2$ orbit basis elements, and hence has norm at most $m+2$.
TB1 bounds the torsion of its cokernel by $(m+2)^A$; the row's
$m$th homology is a subgroup of that cokernel, so it has the same
upper bound on torsion.  This is the required bound on
$T(E^2_{m,0})$.

For $r\ge2$ the homological differential has bidegree $(-r,r-1)$.
There is no incoming differential at $(m,0)$, since its potential
source has second index $1-r<0$.  Thus $E^\infty_{m,0}$ is a
subgroup of $E^2_{m,0}$, and the torsion bound persists.  The other
graded pieces on total degree $m$ are finite and bounded by their
$E^1$ groups.  In an exact sequence $0\to U\to V\to W\to0$ with
finite torsion subgroups, the map $T(V)\to T(W)$ has kernel in
$T(U)$, so $|T(V)|\le |T(U)|\,|T(W)|$.  Applying this to the
finite abutment filtration gives
\[
 |T(H_m(G,\mathbb Z))|
 \le (m+2)^A s^{A\sum_{q=1}^m s^q}
 \le (m+2)^A s^{m A s^m}.
\]
No stabilizer order has been inverted.  Thus this argument retains
small-prime torsion as well as all other torsion.
\end{proof}

Fix henceforth a totally complex number field $L$, with ring of
integers $\mathcal O$, degree $n=[L:\mathbb Q]$ and absolute
discriminant $D\ge1$.  Let $N\ge1$, $d=nN$ and
$\Gamma=\operatorname{GL}_N(\mathcal O)$.  Choose one embedding
$\sigma_w:L\hookrightarrow\mathbb C$ for each conjugate pair,
and put
\[
 V=\prod_w\mathbb C^N,\qquad
 E=\prod_w\operatorname{Herm}_N(\mathbb C),\qquad
 \Lambda=\mathcal O^N\subset V.
 \tag{TB.5}
\]
Thus $\dim_{\mathbb R}V=d$ and
$\dim_{\mathbb R}E=nN^2/2\le dN$.  For $h=(h_w)\in E$ set
\[
 q_h(y)=\sum_w \sigma_w(y)^*h_w\sigma_w(y),\qquad
 \mathcal R=\{h:q_h(y)\ge1\text{ for }0\ne y\in\Lambda\}.
 \tag{TB.6}
\]
Let $\mathcal P$ be the cone of tuples of positive semidefinite
matrices and $\mathcal P^\circ$ its positive definite interior.
Let $U\subset E^*$ be the cone of functionals strictly positive
on $\mathcal P\setminus\{0\}$.  It is open, nonempty and convex:
the sum of the matrix traces belongs to it, and positivity has a
positive minimum on the compact unit section of $\mathcal P$.

\begin{proposition}[TB3: the locally finite Hermitian polyhedron]
\label{prop:tb-hermitian-geometry}
The set $\mathcal R$ is nonempty, closed, convex and contained in
$\mathcal P^\circ$.  Its inequalities are locally finite there.
A point $v\in\mathcal R$ is an extreme point precisely when its
active functionals $h\mapsto q_h(y)$, with $q_v(y)=1$, span $E^*$.
For every $u\in U$, the minimum of $u$ on $\mathcal R$ is attained
at such a vertex.  Only finitely many vertices minimize a functional
in any fixed compact subset of $U$.
\end{proposition}

\begin{proof}
The embedding of $\mathcal O$ is a full lattice; its covolume is
$2^{-n/2}\sqrt D$.  This follows by taking an integral basis,
writing its complex embeddings as real and imaginary parts, and
comparing the square of the embedding determinant with the trace
discriminant.  The product embedding makes $\Lambda$ a full lattice
in $V$.  A sufficiently large scalar multiple of the identity form
therefore belongs to $\mathcal R$.  Closedness and convexity follow
from its defining half-spaces.

If $q_h$ is negative in a real direction, rational approximation in
a lattice basis and multiplication by a common denominator give a
nonzero lattice vector on which it is negative.  Thus $h\in\mathcal R$
is positive semidefinite.  It cannot be singular.  In fact, a
symmetric convex box which is arbitrarily long in its nullspace
and sufficiently thin in the complementary directions has
$q_h<1$ everywhere and arbitrarily large volume.  Minkowski's
convex-body argument then supplies a nonzero lattice point in it,
a contradiction.  The argument used here has the following short
form.  A measurable set of volume greater than the covolume of a
lattice has two distinct points with the same image in a fundamental
parallelepiped, by adding the measures of its translates in that
parallelepiped.  Apply this to half a symmetric convex body of
volume greater than $2^d$ times the covolume.  The difference of
the two points lies in the original body and is a nonzero lattice
point.  This proves the needed version of Minkowski's theorem.

Near a positive definite $h$, all forms have least eigenvalue at
least some $\delta>0$.  A defining inequality that can fail or become
an equality there must therefore have $\|y\|^2\le1/\delta$.
There are only finitely many such lattice vectors.  This proves
local finiteness.  If the active functionals at a point fail to
span, a nonzero direction annihilates all of them; sufficiently
small displacements in either sign remain feasible by local
finiteness.  The point is not extreme.  Conversely, if the active
functionals span and the point is the midpoint of two feasible
points, all its active equalities hold at both endpoints.  Their
difference is annihilated by a spanning set and hence is zero.

For $u\in U$ there is $c>0$ such that $u(h)\ge c\|h\|$ on
$\mathcal P$.  A nonempty sublevel of $u$ on $\mathcal R$ is
therefore compact.  Its minimum face is a nonempty compact convex
set.  Successively minimizing coordinates in a basis produces an
extreme point of that face, which is also extreme in $\mathcal R$.

Finally let $K\subset U$ be compact and fix $h_0\in\mathcal R$.
Uniformly on $u\in K$, one has $u(h)\ge c\|h\|$ on $\mathcal P$
and $u(h_0)\le M$, for some $c>0$ and $M<\infty$.  All minimizing
vertices lie in the compact set
$\mathcal R\cap\{\|h\|\le M/c\}$, which is contained in
$\mathcal P^\circ$.  On this set the least eigenvalue has a positive
minimum.  Thus all active lattice vectors belong to one finite set.
Each vertex is the unique solution to a full-rank selection of the
associated active equalities, so there are only finitely many.
\end{proof}

Let $\mathcal V$ be the vertex set from TB3.  Let $X_p$ consist of
ordered $(p+1)$-tuples of these vertices which simultaneously
minimize some $u\in U$.  Repeated entries are permitted.  These
are exactly the source sets \texttt{Sx}$(L,N,p)$: the forms,
active-constraint definition of a vertex, positive dual cone,
and minimization condition are the same.  Deleting one entry is
the face map.  The action is
\[
       (g\cdot h)_w=(\sigma_w(g)^{-1})^*h_w\sigma_w(g)^{-1}.
       \tag{TB.7}
\]
It preserves $\mathcal R$, vertices and common minimizing faces.

\begin{proposition}[TB4: acyclicity and a common minimal basis]
\label{prop:tb-nerve}
The augmented ordered-tuple complex (TB.2) for these $X_p$ is exact.
For every tuple in $X_p$, its common minimal vectors
\[
 \{y\in\Lambda:q_{v_j}(y)=1\text{ for all }j\}
 \tag{TB.8}
\]
span $L^N$ over $L$.
\end{proposition}

\begin{proof}
For each vertex $v$ let
$C_v=\{u\in U:u(v)\le u(h)\text{ for every }h\in\mathcal R\}$.
These are relatively closed convex subsets of $U$.  They cover $U$
and form a locally finite family by TB3.  A set of vertices is a
simplex of their nerve precisely when its members have a common
minimizing functional.

Here is a finite argument for the augmented homological acyclicity
of this nerve.  A finite cycle uses finitely many simplices.  Choose
a point in the intersection defining each of those simplices,
and let $K\subset U$ be a compact convex set containing these
finitely many points.  Only finitely many $C_v$ meet $K$.  Their
intersections $A_v=C_v\cap K$ are compact convex sets covering
$K$, and the original cycle belongs to their finite nerve.  For
all sufficiently small $\eta>0$, the open $\eta$-neighborhoods
of these sets have the same nerve.  Indeed, an empty intersection
of finitely many compact sets has a positive minimum of the
maximum distance to those sets; there are only finitely many
subfamilies to check.  The union of these neighborhoods is exactly
the $\eta$-neighborhood of $K$, which is convex, and every nonempty
intersection is convex.

The nerve of this finite open convex cover is homologically
acyclic.  One can see the assertion directly from the double
complex of singular chains on the intersections.  In one direction
the straight-line contractions of the nonempty intersections leave
exactly the augmented nerve complex.  In the other direction the
augmented cover complex is exact on chains small enough to lie in
one cover member: for each such singular simplex, its possible
indices form a nonempty simplex, contracted by adjoining any one
of its indices.  Repeated barycentric subdivision makes any finite
collection of chains this small, and the subdivision prism is a
chain homotopy.  The resulting homology is the homology of the
convex union, which vanishes after augmentation.  Thus the chosen
cycle bounds already in the finite nerve, and hence in the whole
nerve.  Nonemptiness follows from TB3.

To pass to ordered tuples, choose any order on the vertex set.
Map an ordered tuple with distinct entries to its increasing order
with the permutation sign, and map tuples with repetitions to
zero.  This is a chain map onto the oriented nerve complex, with
the increasing-tuples inclusion as a section.  On all tuples with
support in a fixed simplex, the augmented complex is contracted
by prepending one chosen vertex.  This proves the section followed
by the projection is homotopic to the identity: inductively in
degree, the difference on a basis tuple, after subtracting the
already constructed homotopy of its boundary, is a cycle supported
on that same simplex; fill it by the prepending contraction.
The ordered complex consequently has the same augmented homology,
and is exact.  The order and the contraction need not be equivariant;
the original ordered complex itself is a permutation complex with
the action (TB.7).

For the spanning assertion let $h$ be the average of the vertices
of the tuple.  Since every value $q_{v_j}(y)$ is at least one,
$q_h(y)=1$ holds exactly for the common minimal vectors (TB.8).
If these vectors did not span $L^N$, a nonzero $L$-linear functional
$a:L^N\to L$ would annihilate them.  At the complex places the
forms $|\sigma_w(a)y|^2$ define a nonzero $k\in\mathcal P$ with
$q_k(y)=0$ on every active constraint at $h$.  Local finiteness
then implies $h-\epsilon k\in\mathcal R$ for all sufficiently
small $\epsilon>0$.  More explicitly, choose a neighborhood with
only finitely many potentially active inequalities; the active
ones do not change, and each other one has a positive margin.
But the common minimizing functional $u$ has $u(k)>0$, and
$u(h-\epsilon k)<u(h)$ contradicts minimality.  This proves the
span assertion and permits a choice of $N$ independent common
minimal vectors in $\mathcal O^N$.
\end{proof}

Put
\[
       I_0=D^N d^d,\qquad R_0=N^N I_0^2,
       \qquad V_0=(2R_0+1)^d.
       \tag{TB.9}
\]
These are the source integers \texttt{I0}, \texttt{RR}, and the
near-vector bound.  They are positive.

\begin{proposition}[TB5: determinant and common-basis index bounds]
\label{prop:tb-determinant}
For $h\in\mathcal R$ one has
\[
                 1\le I_0\prod_w\det h_w.
                 \tag{TB.10}
\]
If $A\in\operatorname{Mat}_N(\mathcal O)$ is nonsingular and all
its columns are minimal for $h$, then
$\Delta=|N_{L/\mathbb Q}(\det A)|$ satisfies
$1\le\Delta\le I_0$.  Moreover the matrices
$H_w=\sigma_w(A)^*h_w\sigma_w(A)$ satisfy
\[
 H_w>0,\qquad (H_w)_{ii}\le1,\qquad
                 I_0^{-1}\le\det H_w\le1.
                 \tag{TB.11}
\]
\end{proposition}

\begin{proof}
Let $G_h$ be the real Gram matrix of $q_h$ in an integral basis of
$\Lambda$.  Its minimum on $\mathbb Z^d\setminus\{0\}$ is at
least one.  The inverse image under $G_h^{1/2}$ of the open cube
$(-d^{-1/2},d^{-1/2})^d$ contains no nonzero lattice point.
The convex-body argument in TB3 therefore bounds its volume by
$2^d$.  This gives $\det G_h\ge d^{-d}$.  The discriminant
calculation used in TB3 and the real determinant of a Hermitian
form give the exact identity
\[
        2^d\det G_h=D^N\left(\prod_w\det h_w\right)^2.
        \tag{TB.12}
\]
Consequently $\prod_w\det h_w\ge\sqrt{2^d/I_0}\ge I_0^{-1}$,
which proves (TB.10).

For a minimal column of $A$, the sum over $w$ of the corresponding
diagonal entries of $H_w$ is one.  These entries are nonnegative,
so each is at most one.  Hadamard's inequality gives
$0<\det H_w\le1$.  The product formula for the determinant is
\[
       \prod_w\det H_w
           =\Delta\prod_w\det h_w.
           \tag{TB.13}
\]
Here $\Delta$ is a positive integer.  Since the left side is at
most one, (TB.10) implies $\Delta\le I_0$.  Since all factors
on the left are at most one, each one is at least their product,
which by (TB.10) and $\Delta\ge1$ is at least $I_0^{-1}$.
This proves (TB.11).
\end{proof}

\begin{proposition}[TB6: a finite left-matrix quotient without normal forms]
\label{prop:tb-row-module}
For every integer $J\ge1$, the nonsingular matrices
$A\in\operatorname{Mat}_N(\mathcal O)$ with
$|N_{L/\mathbb Q}(\det A)|\le J$ have at most $J^{1+dN}$
orbits under left multiplication by $\Gamma$.
\end{proposition}

\begin{proof}
Let $M_A\subset\mathcal O^N$ be the module generated by the rows
of $A$.  Its index as an abelian subgroup is
$I=|N_{L/\mathbb Q}(\det A)|$.  This follows either by the real
embedding determinant, or by the determinant of the induced
$\mathbb Q$-linear map on $L^N$ in an integral basis.
Since the finite quotient has order $I$, it is annihilated by $I$,
so $I\mathcal O^N\subset M_A$.  The image of $M_A$ in
$(\mathcal O/I\mathcal O)^N$ is generated over $\mathcal O/I\mathcal O$
by its $N$ row residues.  There are at most
\[
             (I^{nN})^N=I^{dN}
\]
possible ordered lists of such residues, and each list determines
its generated submodule and hence its inverse image $M_A$.

Two nonsingular matrices have the same row module precisely when
their rows are two $\mathcal O$-bases of that module, and their
basis-change matrix belongs to $\operatorname{GL}_N(\mathcal O)$.
Thus they are in the same left orbit.  Summing the preceding bound
over $1\le I\le J$ gives
$\sum_{I=1}^J I^{dN}\le J^{1+dN}$.  No principal-ideal hypothesis
on $\mathcal O$ was used; the row modules here are free because
their displayed rows are independent over $L$.
\end{proof}

\begin{proposition}[TB7: near-vector and vertex counts]
\label{prop:tb-near-count}
Fix a nonsingular integral $A$ of determinant norm
$\Delta\le I_0$.  The union of all minimal vectors of all
$h\in\mathcal R$ for which the columns of $A$ are minimal has
cardinality at most $V_0$.  The number of such forms which are
vertices is at most $(V_0+1)^{dN}$.
\end{proposition}

\begin{proof}
The column lattice $A\mathcal O^N$ has index $\Delta$ in
$\mathcal O^N$.  Thus $\Delta\mathcal O^N\subset A\mathcal O^N$,
and there is an integral matrix $B$ with
$AB=BA=\Delta I$.  It is nonsingular.  Let $y$ be a vector in
the stated union and choose an associated form $h$.
At a complex place put $c=\sigma_w(A)^{-1}\sigma_w(y)$.
Then $c^*H_wc\le q_h(y)=1$.  The inverse-matrix Cauchy inequality
gives
\[
        |c_i|^2\le (H_w^{-1})_{ii}\,(c^*H_wc).
        \tag{TB.14}
\]
Indeed this is ordinary Cauchy--Schwarz for the inner product with
matrix $H_w$, applied to $c$ and $H_w^{-1}e_i$.
The diagonal cofactor of $H_w$ is at most one by Hadamard's
inequality and (TB.11), including the empty cofactor when $N=1$.
Therefore $(H_w^{-1})_{ii}\le I_0$ and
\[
       |\sigma_w((By)_i)|=\Delta |c_i|
                \le I_0^{3/2}\le R_0.
                \tag{TB.15}
\]
The same bound holds at conjugate embeddings.

There are at most $(2R_0+1)^n$ algebraic integers all of whose
conjugates have absolute value at most $R_0$.  To prove this,
reduce them modulo $Q\mathcal O$ where $Q=2R_0+1$.  If two
different such integers have difference $Qz$, the nonzero
integer $N_{L/\mathbb Q}(z)$ has absolute value at least one.
Some conjugate of $z$ consequently has absolute value at least
one, whereas the corresponding conjugate of the difference has
absolute value at most $2R_0<Q$, a contradiction.  Reduction is
injective, and $|\mathcal O/Q\mathcal O|=Q^n$ because
$\mathcal O$ is free of rank $n$ over $\mathbb Z$.
Applying this coordinatewise to $By$ and using injectivity of $B$
gives at most $Q^{nN}=V_0$ choices for $y$.

At a vertex the active functionals span $E^*$ by TB3.  Choose
at most $\dim E\le dN$ minimal vectors whose functionals form
a basis.  Their equations $q_h(y)=1$ determine the vertex
uniquely.  Each chosen vector lies in the preceding finite union.
Choose one ordered list for each vertex and pad to length $dN$
with a dummy symbol.  Two vertices with the same list agree by
the full-rank equalities.  Hence the number of vertices is at
most $(V_0+1)^{dN}$, as asserted.
\end{proof}

\begin{proposition}[TB8: all orbit and stabilizer budgets]
\label{prop:tb-orbits-stabilizers}
For every $p\ge0$, the $\Gamma$-orbits of $X_p$ are finite in
number and satisfy
\[
 \begin{split}
 |\Gamma\backslash X_p|
 &\le A_p:=I_0^{1+dN}(V_0+1)^{dN(p+1)}\\
 &\le (2Dd)^{12(p+1)d^3N}.
 \end{split}
 \tag{TB.16}
\]
Every tuple stabilizer is finite of order at most
\[
                         s=3^{dN}.
                         \tag{TB.17}
\]
\end{proposition}

\begin{proof}
TB4 supplies $N$ independent common minimal vectors for any tuple;
make them the columns of $A$.  TB5 gives $\Delta\le I_0$.
By TB6, a left translate $gA$ can be chosen from a list of at
most $I_0^{1+dN}$ matrices.  Under (TB.7) the minimal vectors
translate by $g$, so every vertex of the translated tuple belongs
to the near-vertex set for this representative.  TB7 gives at
most $(V_0+1)^{dN(p+1)}$ such ordered tuples.  This proves the
first bound, without requiring unique choices of common basis
or orbit representative.

For the numerical bound put $B_0=2Dd\ge2$.  Since $N\le d$,
$D,d,N\le B_0$, and
\[
 I_0\le B_0^{2d},\qquad R_0\le B_0^{5d},\qquad
 2R_0+1\le B_0^{5d+2},\qquad
 V_0+1\le B_0^{(5d+2)d+1}.
\]
The exponent in the resulting upper bound on $A_p$ is
\[
 2d(1+dN)+\bigl((5d+2)d+1\bigr)dN(p+1)
                    \le12(p+1)d^3N.
\]
For an explicit check, the terms $2d$, $2d^2N$,
$5d^3N(p+1)$, $2d^2N(p+1)$ and $dN(p+1)$ are bounded
respectively by $2,2,5,2,1$ times $(p+1)d^3N$.

A stabilizer of a tuple lies in the stabilizer of its first
positive definite form.  Under all complex embeddings that
stabilizer lies in a product of compact unitary groups.  Its
matrix entries lie in the lattice $\operatorname{Mat}_N(\mathcal O)$,
so it is finite.  A finite subgroup of $\Gamma$ injects under
reduction modulo $3\mathcal O$.  Here is the needed torsion-free
kernel argument.  If a nonidentity element of that kernel has
finite order, a power has prime order $q$; write it
$I+3^r C$, where $r\ge1$ is maximal and some integral entry of
$C$ is not divisible by $3$.  For $q\ne3$, the binomial expansion
of its $q$th power, reduced modulo $3^{r+1}$, has nonzero term
$q3^r C$.  For $q=3$, reduction modulo $3^{r+2}$ has nonzero
term $3^{r+1}C$; the remaining terms have factors
$3^{2r+1}$ and $3^{3r}$ and vanish in that quotient.  Both
contradict the claimed order.  The maximal $r$ exists because
$\mathcal O$ is a finite free abelian group.
Finally
$|\operatorname{Mat}_N(\mathcal O/3\mathcal O)|
=3^{nN^2}=3^{dN}$, which bounds the size of the finite subgroup
and proves (TB.17).
\end{proof}

\begin{theorem}[TB9: the literal arithmetic torsion bound]
\label{thm:tb-literal}
For every totally complex number field $L$ and every integer
$e\ge1$, set $m=2e$, $N=4e+4$ and $d=nN$.
The torsion subgroup of $H_m(\operatorname{GL}_N(\mathcal O),\mathbb Z)$
is finite, and
\[
 \log|T(H_m(\operatorname{GL}_N(\mathcal O),\mathbb Z))|
     \le \exp\!\bigl(49\,m d^3N\log(2Dd)\bigr).
 \tag{TB.18}
\]
This is exactly \texttt{TorsBound}$(49)$.
\end{theorem}

\begin{proof}
Use the exact complex of TB4 in TB2.  The integers $A_p$ in
TB8 increase with $p$, so take $A=A_m$ and $s=3^{dN}$.
Put $X=m d^3N\log(2Dd)$.  Since $m\ge2$, (TB.16) gives
$A\le e^{24X}$.  Also $X\ge1$,
$\log(m+2)\le X$, and
$\alpha=m\log s=m dN\log3\le X$; these follow respectively
from $d\ge2$, $N\ge1$, $\log(m+2)\le m$ for $m\ge2$,
and $d^2\log(2Dd)\ge\log3$.
Taking logarithms in (TB.3) yields
\[
 \log|T(H_m)|
    \le A\bigl(\log(m+2)+m s^m\log s\bigr)
    \le e^{24X}(X+Xe^X)
    \le e^{27X}\le e^{49X}.
\]
In the penultimate step use $X\le e^X$ and
$2e^{2X}\le e^{3X}$ for $X\ge1$.
The source's \texttt{HGL} is this group homology with the
trivial integral representation, and \texttt{torsHGL} is its
full torsion subgroup.  Its \texttt{tBound} is the right side
of (TB.18), with $n$ the original degree and $D$ the original
absolute discriminant.  The finiteness and inequality therefore
supply the literal predicate with all its field quantifiers.
\end{proof}

\begin{theorem}[TB10: an explicit power supplier for LemmaK]
\label{thm:tb-power}
For all the fields and integers in TB9,
\[
 \log|T(H_{2e}(\operatorname{GL}_{4e+4}(\mathcal O),\mathbb Z))|
        \le \exp\!\bigl((2en)^{22}\log(2D)\bigr).
        \tag{TB.19}
\]
Thus the source's \texttt{TorsBoundPow}$(22)$ holds, with the
constant chosen before the number field and $e$.
\end{theorem}

\begin{proof}
Put $R=2en\ge2$.  Then $m\le R$,
$N=4e+4\le8e\le4R$, and $d=nN\le4R$.
For $R\ge2$,
$\log(4R)\le2R\log2$; the inequality holds at $R=2$ and the
derivative of the right side minus the left side is positive
thereafter.  Consequently
\[
 \log(2Dd)\le\log(2D)+\log(4R)
           \le3R\log(2D).
\]
It follows that
\[
 \begin{split}
 49m d^3N\log(2Dd)
 &\le 49R(4R)^3(4R)(3R)\log(2D)\\
 &=37632R^6\log(2D)
 \le R^{22}\log(2D),
 \end{split}
\]
since $R^{16}\ge2^{16}>37632$.
Insert this into TB9.  The right side is exactly
\texttt{kBound}$(22,e,n,D)$, and the same torsion subgroup is
finite by TB9.  No field extension, replacement discriminant,
or excluded torsion prime is involved.
\end{proof}

\paragraph{Closed endpoint.}
TB10 supplies the complete torsion premise in the K-theoretic
reduction.  That reduction still needs its actual comparison
maps and the finite local/global cohomology argument; these
are supplied separately in the K-theory chapter.  The present
chapter alone makes no claim about a common arithmetic family,
the analytic gap, RH or Dirichlet GRH.
